\documentclass[conference]{IEEEtran}
\IEEEoverridecommandlockouts
%Template version as of 6/27/2024

\usepackage{cite}
\usepackage{amsmath,amssymb,amsfonts}
\usepackage{algorithmic}
\usepackage{graphicx}
\usepackage{textcomp}
\usepackage{xcolor}
\usepackage[utf8]{inputenc}
\usepackage[T1]{fontenc}
\usepackage{booktabs}

\def\BibTeX{{\rm B\kern-.05em{\sc i\kern-.025em b}\kern-.08em
    T\kern-.1667em\lower.7ex\hbox{E}\kern-.125emX}}

\begin{document}

\title{Detección Automática de Fracturas Óseas en Radiografías Mediante Modelos de Visión por Computadora Basados en Deep Learning}

\author{
\IEEEauthorblockN{Lucia T. Capon Paul}
\IEEEauthorblockA{\textit{CEIA -- FIUBA} \\
Buenos Aires, Argentina}
\and
\IEEEauthorblockN{Cesar Orellana}
\IEEEauthorblockA{\textit{CEIA -- FIUBA} \\
Buenos Aires, Argentina}
\and
\IEEEauthorblockN{Leandro Britez}
\IEEEauthorblockA{\textit{CEIA -- FIUBA} \\
Buenos Aires, Argentina}
}

\maketitle

% ─── ABSTRACT ────────────────────────────────────────────────────────────────
\begin{abstract}
% TODO: completar tras finalizar los experimentos (~150 palabras)
% Incluir: problema, enfoque, modelos comparados (YOLOv11 vs RT-DETR),
%          dataset usado, métricas clave obtenidas, y principal conclusión.
La detección automatizada de fracturas óseas en imágenes radiográficas representa
un desafío clínico de alto impacto. En este trabajo se presenta el desarrollo de un
prototipo basado en modelos de detección de objetos preentrenados para la
identificación y localización de fracturas en radiografías.
Se comparan dos arquitecturas: YOLOv11, representativa del paradigma
anchor-free de una etapa, y RT-DETR, un detector basado en transformers de
tiempo real. Los modelos son ajustados (\textit{fine-tuned}) sobre el dataset
público \textit{Bone Fracture Detection} de Roboflow Universe.
La evaluación se realiza en términos de mAP@0.5, mAP@0.5:0.95, F1-score y
latencia de inferencia. [COMPLETAR CON RESULTADOS].
\end{abstract}

\begin{IEEEkeywords}
detección de objetos, fracturas óseas, YOLOv11, RT-DETR, deep learning,
visión por computadora, radiografías.
\end{IEEEkeywords}

% ─── I. INTRODUCCIÓN ─────────────────────────────────────────────────────────
\section{Introducción}

La interpretación de imágenes radiográficas es una tarea clínica fundamental en
el diagnóstico de traumatismos y patologías óseas. La identificación visual de
fracturas requiere experiencia especializada y puede verse afectada por factores
como la carga de trabajo y la fatiga del profesional \cite{b1}.

Los sistemas de visión por computadora basados en \textit{deep learning} han
demostrado capacidad para asistir en tareas de detección y localización de
patologías en imágenes médicas \cite{b2}. En particular, los modelos de
detección de objetos permiten no solo determinar la presencia de una lesión,
sino también localizar su posición mediante cuadros delimitadores
(\textit{bounding boxes}).

% TODO: agregar caso real de aplicación similar (no mayor a 8 años)

El objetivo de este trabajo es desarrollar y evaluar un prototipo funcional
capaz de detectar automáticamente fracturas óseas en radiografías, comparando
el desempeño de dos arquitecturas modernas: YOLOv11 \cite{b3} y RT-DETR \cite{b4}.

% ─── II. METODOLOGÍA ─────────────────────────────────────────────────────────
\section{Metodología, Diseño y Desarrollo}

\subsection{Dataset}

Se utilizó el dataset \textit{Bone Fracture Detection} disponible públicamente
en Roboflow Universe \cite{b5}. El dataset contiene imágenes de radiografías con
anotaciones en formato YOLO (\textit{bounding boxes}). Las particiones utilizadas
son:

\begin{table}[htbp]
\caption{Distribución del dataset por split}
\begin{center}
\begin{tabular}{lcc}
\toprule
\textbf{Split} & \textbf{Imágenes} & \textbf{Anotaciones} \\
\midrule
Train & -- & -- \\  % TODO: completar con valores reales
Valid & -- & -- \\
Test  & -- & -- \\
\bottomrule
\end{tabular}
\end{center}
\end{table}

\subsection{Preprocesamiento y Data Augmentation}

Todas las imágenes fueron redimensionadas a $640 \times 640$ píxeles. Se
aplicaron las siguientes técnicas de \textit{data augmentation} durante
el entrenamiento: variaciones de matiz y saturación (HSV), flip horizontal,
traslación aleatoria y escalado. Para RT-DETR se desactivó el \textit{mosaic
augmentation} por ser incompatible con la arquitectura \textit{transformer}.

\subsection{Arquitecturas}

\textbf{YOLOv11} \cite{b3}: detector \textit{anchor-free} de una sola etapa.
Se utilizó la variante \texttt{yolo11m.pt} preentrenada en COCO.

\textbf{RT-DETR} \cite{b4}: detector basado en transformers diseñado para
inferencia en tiempo real. Se utilizó la variante \texttt{rtdetr-l.pt}
preentrenada en COCO.

\subsection{Entrenamiento}

Ambos modelos fueron ajustados durante 100 épocas con optimizador
\textit{AdamW}. Los hiperparámetros completos se documentan en el repositorio
de GitHub en los archivos \texttt{configs/yolov11.yaml} y
\texttt{configs/model2.yaml}.

% ─── III. RESULTADOS ─────────────────────────────────────────────────────────
\section{Resultados Experimentales y Análisis}

% TODO: completar esta sección tras correr los experimentos.

\begin{table}[htbp]
\caption{Comparación de métricas en el conjunto de test}
\begin{center}
\begin{tabular}{lcccc}
\toprule
\textbf{Modelo} & \textbf{mAP@0.5} & \textbf{mAP@0.5:0.95} & \textbf{F1} & \textbf{Latencia (ms)} \\
\midrule
YOLOv11  & -- & -- & -- & -- \\
RT-DETR  & -- & -- & -- & -- \\
\bottomrule
\end{tabular}
\end{center}
\end{table}

% Agregar figuras de curvas PR, ejemplos de detección, análisis de errores.

% ─── IV. CONCLUSIONES ────────────────────────────────────────────────────────
\section{Conclusiones}

% TODO: al menos 5 conclusiones (requerimiento de la cátedra)

\begin{enumerate}
    \item ...
    \item ...
    \item ...
    \item ...
    \item ...
\end{enumerate}

% ─── BIBLIOGRAFÍA ────────────────────────────────────────────────────────────
\begin{thebibliography}{00}

\bibitem{b1}
% TODO: completar con paper de Papers/ relevante
Author, ``Title,'' Journal, vol. X, pp. Y--Z, year.

\bibitem{b2}
% TODO: completar con paper de Papers/
Author, ``Title,'' in Proc. Conference, year, pp. X--Y.

\bibitem{b3}
Ultralytics, ``YOLO11: State-of-the-Art Object Detection,'' 2024.
[Online]. Available: https://github.com/ultralytics/ultralytics

\bibitem{b4}
L. Leng et al., ``DETRs Beat YOLOs on Real-time Object Detection,''
arXiv:2304.08069, 2023.

\bibitem{b5}
Roboflow Universe, ``Bone Fracture Detection Dataset,'' 2024.
[Online]. Available: https://universe.roboflow.com/veda/bone-fracture-detection-daoon

\end{thebibliography}

\end{document}
